\section{Additional Experimental Details}
\label{appx:experiment-details}

\subsection{Retrieval Details}
\label{appx:retrieval}

\textbf{Sparse retrieval.}
During retrieval we make a slight augmentation to the documents by pre-pended files' contents with their file paths to better enable retrieval based on filenames that may be mentioned directly in the issue.

\textbf{Oracle retrieval.}
Oracle retrieval file paths are simply extracted directly from the reference solution's patch file excluding test files.

\subsection{Inference Settings}
Since generations are relatively expensive, we only generate a single patch file per instance. Following precedent in code generation for evaluation in Pass@$1$~\citep{chen2021evaluating,rozière2023code}, we simply use greedy decoding for all models.

\subsection{Prompt Template Example}
\label{appx:experimental-setup:prompt-template-example}
Models are prompted with the following general template with slight variations depending on the model used.
\begin{lstlisting}
You will be provided with a partial code base and an issue statement explaining a problem to resolve.
<issue>
{ISSUE TEXT}
</issue>

<code>
[start of README.md]
{README.md text}
[end of README.md]
[start of file_1.py]
{file_1.py text}
[end of file_1.py]
...
</code>

Here is an example of a patch file. It consists of changes to the code base. It specifies the file names, the line numbers of each change, and the removed and added lines. A single patch file can contain changes to multiple files.

<patch>
--- a/file.py
+++ b/file.py
@@ -1,27 +1,35 @@
 def euclidean(a, b):
-    while b:
-        a, b = b, a % b
-    return a
+    if b == 0:
+        return a
+    return euclidean(b, a % b)
 
 
 def bresenham(x0, y0, x1, y1):
     points = []
     dx = abs(x1 - x0)
     dy = abs(y1 - y0)
-    sx = 1 if x0 < x1 else -1
-    sy = 1 if y0 < y1 else -1
-    err = dx - dy
+    x, y = x0, y0
+    sx = -1 if x0 > x1 else 1
+    sy = -1 if y0 > y1 else 1
 
-    while True:
-        points.append((x0, y0))
-        if x0 == x1 and y0 == y1:
-            break
-        e2 = 2 * err
-        if e2 > -dy:
+    if dx > dy:
+        err = dx / 2.0
+        while x != x1:
+            points.append((x, y))
             err -= dy
-            x0 += sx
-        if e2 < dx:
-            err += dx
-            y0 += sy
+            if err < 0:
+                y += sy
+                err += dx
+            x += sx
+    else:
+        err = dy / 2.0
+        while y != y1:
+            points.append((x, y))
+            err -= dx
+            if err < 0:
+                x += sx
+                err += dy
+            y += sy
 
+    points.append((x, y))
     return points
</patch>

I need you to solve the provded issue by generating a single patch file that I can apply directly to this repository using git apply. Please respond with a single patch file in the format shown above.

Respond below:
\end{lstlisting}

Experiments using slightly more or fewer lines of instructions or examples seemed to not affect overall performance substantially, except for the findings of experiments stated in Section~\ref{sec:results}.
